% Fragmento integrable. Preambulo anfitrion: amsmath, amssymb, longtable, hyperref.
\section{Classification et prolongement du sélecteur des premières racines}
\label{hmtfg:fragmento:clasificacion-y-continuacion-del-selector-de-primeras-raices}

\subsection{Objet et provenance}
\label{hmtfg:clasificacion:1}

Cette annexe conserve la famille issue de la lecture dodécaphasique HMT :

\[
u_0(x)=\frac1{12}\sum_{m=1}^{12}\left(1-\cos\frac{2(3m-1)x}{\sqrt{899}}\right),
\qquad f_\lambda(x)=1-\lambda u_0(x).
\]

La famille provient des antécédents operatoriels et analytiques exposés dans ce développement. L’existence et la compacité de ses réalisations de l’itinéraire infini ont été démontrées dans les sous-sections consacrées à la réalisation du mot et à la conservation de ses signes.

L’argument suivant est propre à ce développement et utilise exclusivement ces propriétés et les opérations de retour de la famille fixée. Il n’utilise comme entrée ni un paramètre d’une autre famille ni la valeur d’une constante universelle. Les feuilles des itinéraires utilisés ici sont les deux branches monotones du lecteur critique ; les autres coordonnées et la mémoire HMT conservent leur résidence dans l’état dont provient le lecteur.

Écrivons
\[
c_j(\lambda)=f_\lambda^j(0),\qquad
\tau_j=(-1)^{v_2(j)},\qquad
W_n=\tau_1\cdots\tau_{2^n-1}.
\]
L’itinéraire d’une orbite critique qui revient pour la première fois à l’origine se termine par le symbole zéro. L’ordre des itinéraires est l’ordre unimodal à maximum : avant la comparaison du symbole d’indice \(j\), le facteur d’orientation est
\[
O_{j-1}(K)=\prod_{i=1}^{j-1}(-K_i).
\]
L’ordre ordinaire des symboles est \(-1<0<1\).


\subsection{Classification avant le mot infini}
\label{hmtfg:clasificacion:2}

\textbf{Théorème de classification.} Soit \(F:I\to I\), avec \(I=[\ell,1]\), \(\ell<0\), une application continue, strictement croissante à gauche de zéro et strictement décroissante à sa droite, dont le maximum satisfait \(F(0)=1\). Supposons que zéro soit de période exacte \(2^n\), avec \(n\ge1\), et que son itinéraire critique \(K\) satisfasse \(K<\tau\). Alors
\[
\boxed{K=W_n0.}
\]
La conclusion porte sur l’itinéraire et permet que des paramètres distincts d’une famille possèdent ce même itinéraire.

\textbf{Démonstration.} Notons \(d_j=F^j(0)\). Pour la période deux, l’itinéraire est \(+0=W_10\). Considérons une période exacte \(p=2^n\ge4\).

On a \(d_2<0\). En effet, \(d_2=0\) donnerait la période deux. Si \(d_2>0\), la monotonie de la branche positive donne
\[
F([d_2,1])=[d_2,d_3]\subset[d_2,1],
\]
et l’orbite critique demeure dans un intervalle strictement positif, incompatible avec le retour à zéro.

Les signes initiaux sont alors \(+,-\). Le facteur d’orientation pour comparer le troisième signe est \(-1\) ; par conséquent, \(K<\tau\) impose \(d_3>0\). Un zéro à cette position produirait la période trois.

On ne peut pas non plus avoir \(d_4<0\). Sous cette hypothèse,
\[
d_2<d_4<0,
\qquad F([d_2,d_4])=[d_3,d_5]\subset(0,1],
\]
car \(d_4=F(d_3)>F(1)=d_2\) et \(d_5=F(d_4)>F(d_2)=d_3\). La branche positive est décroissante, de sorte que
\[
F^2([d_2,d_4])=[d_6,d_4]\subset[d_2,d_4].
\]
Ici \(d_6\ge d_2\), car \(d_5\le1\), et \(d_6<d_4\), car \(d_5>d_3\). Cet intervalle strictement négatif piège les retours pairs de l’orbite critique et exclut le retour à zéro. Donc \(d_4\ge0\). Si \(p=4\), on obtient précisément \(+-+0=W_20\).

Supposons maintenant \(p\ge8\). Alors \(d_4>0\). Les facteurs d’orientation pour les cinquième et sixième signes sont, respectivement, \(-1\) et \(+1\). Comme \(\tau_5=+1\) et \(\tau_6=-1\), l’inégalité \(K<\tau\) impose
\begin{equation}
\operatorname{sign}(d_1,\ldots,d_6)=(+,-,+,+,+,-).
\label{hmtfg:clasificacion:eq:1}
\end{equation}
Les signes nuls à ces positions sont exclus par la période exacte.

Nous construisons l’intervalle de retour
\[
J=[d_2,d_4].
\]
L’inégalité \(F(d_3)=d_4>0>d_6=F(d_5)\), jointe à \(d_3,d_5>0\), implique \(d_3<d_5\) ; donc
\[
F(J)=[d_3,1].
\]
De plus, \(d_4<d_3\). Pour le vérifier, l’image par \(F\) de l’intervalle d’extrémités \(d_3,d_4\) est strictement positive, puisque les images de ses extrémités sont \(d_4,d_5>0\). Sur cet intervalle, \(F^2\) est strictement croissante. Ses valeurs aux extrémités satisfont
\[
F^2(d_3)=d_5>0>d_6=F^2(d_4),
\]
ce qui impose \(d_3>d_4\). Par conséquent,
\begin{equation}
J\cap F(J)=\varnothing,
\qquad F^2(J)=F([d_3,1])=[d_2,d_4]=J.
\label{hmtfg:clasificacion:eq:2}
\end{equation}

L’application \(F^2|_J\) possède un unique minimum en zéro. La conjugaison par \(h(x)=d_2x\), qui inverse l’orientation, définit
\begin{equation}
G=h^{-1}\circ F^2\circ h:
\left[\frac{d_4}{d_2},1\right]\longrightarrow
\left[\frac{d_4}{d_2},1\right].
\label{hmtfg:clasificacion:eq:3}
\end{equation}
Cette application satisfait les mêmes hypothèses de maximum que \(F\), avec \(G(0)=1\), et son point critique a pour période exacte \(p/2\). Son itinéraire est
\begin{equation}
K'_j=-K_{2j}.
\label{hmtfg:clasificacion:eq:4}
\end{equation}
Tous les symboles d’indice impair de \(K\) sont positifs, car les positions paires appartiennent à \(J\) et les positions impaires à \(F(J)\subset(0,1]\).

La transformation \eqref{hmtfg:clasificacion:eq:4} conserve l’ordre de comparaison avec \(\tau\). En effet, \(\tau_{2j-1}=+1\) et \(\tau_{2j}=-\tau_j\). Si \(q\) est la première position où \(K'\) et \(\tau\) diffèrent, la première position de différence entre \(K\) et \(\tau\) est \(2q\), et
\begin{equation}
O_{2q-1}(K)=-O_{q-1}(K'),
\qquad K_{2q}-\tau_{2q}=-(K'_q-\tau_q).
\label{hmtfg:clasificacion:eq:5}
\end{equation}
Le produit qui détermine l’ordre est identique. Les identités restent valides lorsque le premier symbole distinct est le zéro final. Ainsi \(K'<\tau\).

La récurrence appliquée à \(G\) donne \(K'=W_{n-1}0\). Les positions impaires positives et \eqref{hmtfg:clasificacion:eq:4} reconstruisent exactement \(K=W_n0\). Le théorème est démontré. \(\square\)

Le théorème démontre également que chaque application classifiée possède les retours restrictifs de doublement antérieurs à son retour superstable, avec les domaines de \eqref{hmtfg:clasificacion:eq:2}–\eqref{hmtfg:clasificacion:eq:3}. Il ne présuppose pas de monotonie du paramètre d’une famille.

\textbf{Normalisation symétrique de la famille fixée.} Lorsque \(F\) est paire et définie sur \([-1,1]\), les inégalités précédentes donnent en outre
\[
F(d_4)=d_5>d_3=F(d_2)=F(-d_2).
\]
Comme les deux arguments \(d_4,-d_2\) sont positifs et que \(F\) décroît sur cette branche,
\begin{equation}
0<d_4<-d_2<1.
\label{hmtfg:clasificacion:eq:5a}
\end{equation}
Donc \(F^2([d_2,-d_2])=[d_2,d_4]\), et la même conjugaison \(h(x)=d_2x\) définit sur tout \([-1,1]\) une application paire à maximum dont l’image est \([d_4/d_2,1]\subset[-1,1]\). Elle coïncide exactement avec l’opérateur \(RF=-F(F(-ax))/a\), \(a=-F(1)\), utilisé dans le corpus. Cette extension symétrique conserve le retour critique, sa période et son itinéraire.


\subsection{Premier paramètre d’itinéraire infini}
\label{hmtfg:clasificacion:3}

Nous travaillons sur l’intervalle compact déjà utilisé dans la démonstration d’existence,
\[
L=\left[\frac1{2u_0(1)},\frac2{u_0(1)}\right].
\]
Soit
\[
\Lambda_\tau=\{\lambda\in L:
\operatorname{sign}c_j(\lambda)=\tau_j\text{ pour tout }j\ge1\}.
\]
D’après les démonstrations d’existence et de séparation des signes du prolongement, cet ensemble est non vide et compact. Il existe donc
\begin{equation}
a=\min\Lambda_\tau.
\label{hmtfg:clasificacion:eq:6}
\end{equation}

\textbf{Lemme de comparaison.} Pour tout \(\lambda\in[1/(2u_0(1)),a)\), on a \(K_\lambda<\tau\).

\textbf{Démonstration.} À l’extrémité gauche, l’itinéraire est \(+^\infty<\tau\). L’argument de séparation de la réalisation du mot infini démontre que les ensembles de comparaison stricte avec \(\tau\) sont ouverts : dans un itinéraire infini distinct de \(\tau\), la première différence finie persiste par continuité ; à un paramètre superstable, les deux mots périodiques voisins et le mot à zéro final se situent du même côté de \(\tau\), par sa maximalité stricte. L’intervalle antérieur à \(a\) ne contient, par définition, aucun paramètre d’itinéraire \(\tau\). Sa connexité maintient la comparaison initiale sur tout cet intervalle. \(\square\)


\subsection{Existence de chaque première nouvelle racine}
\label{hmtfg:clasificacion:4}

Nous conservons le sélecteur original. La première racine est
\[
\lambda_1=\frac1{u_0(1)},
\]
et, pour \(n\ge2\), \(\lambda_n\) est la première racine de \(c_{2^n}\) située à droite de \(\lambda_{n-1}\) qui satisfait
\begin{equation}
c_{2^j}(\lambda_n)\ne0\quad(1\le j<n).
\label{hmtfg:clasificacion:eq:7}
\end{equation}
Comme toute première période divisant \(2^n\) est une puissance de deux, \eqref{hmtfg:clasificacion:eq:7} équivaut à exiger une période critique exacte \(2^n\).

\textbf{Lemme d’existence.} Toutes les racines de ce sélecteur existent et satisfont
\begin{equation}
\lambda_1<\lambda_2<\cdots<\lambda_n<a.
\label{hmtfg:clasificacion:eq:8}
\end{equation}

\textbf{Démonstration.} Le signe \(c_2(a)<0\) implique \(\lambda_1<a\). Supposons \(\lambda_{n-1}<a\) construite, et posons \(p=2^{n-1}\). La fonction analytique \(c_p(\lambda)\) n’est pas identiquement nulle, puisque \(c_p(a)\ne0\). Son ensemble de zéros dans \([\lambda_{n-1},a]\) est fini, non vide et ne contient pas \(a\). Soit \(z<a\) son dernier zéro.

Sur \((z,a]\), le signe de \(c_p\) est constant et égal à \(s=\tau_p\). Pour \(g_\lambda=f_\lambda^p\), on a
\[
g_\lambda(0)=c_p(\lambda),\qquad g'_\lambda(0)=0.
\]
Une borne locale uniforme de la dérivée seconde donne
\begin{equation}
c_{2p}(\lambda)
=g_\lambda(c_p(\lambda))
=c_p(\lambda)+O(c_p(\lambda)^2).
\label{hmtfg:clasificacion:eq:9}
\end{equation}
Lorsque \(\lambda\) tend vers \(z\) par la droite, le terme du premier ordre domine ; donc \(c_{2p}(\lambda)\) a pour signe \(s\). En \(a\), son signe est
\[
\tau_{2p}=-\tau_p=-s.
\]
Le théorème des valeurs intermédiaires fournit un zéro \(\beta\in(z,a)\) de \(c_{2p}\). Comme \(c_p\ne0\) sur cet intervalle, la période critique en \(\beta\) est exactement \(2p\).

L’analyticité et \(c_{2p}(a)\ne0\) impliquent que les zéros de \(c_{2p}\) dans \([\lambda_{n-1},a]\) sont en nombre fini. Parmi ses zéros de période exacte \(2p\) situés à droite de \(\lambda_{n-1}\), il en existe donc un premier. C’est le sélecteur original \(\lambda_n\), et il satisfait \(\lambda_n\le\beta<a\). \(\square\)

L’utilisation du dernier zéro auxiliaire \(z\) appartient uniquement à la preuve d’existence. Le paramètre sélectionné reste la \textbf{première nouvelle racine} \(\lambda_n\).


\subsection{Convergence du sélecteur global vers le premier itinéraire infini}
\label{hmtfg:clasificacion:5}

\textbf{Théorème de prolongement global.} Pour la famille HMT fixée et le sélecteur original, tous les itinéraires sélectionnés sont canoniques et
\begin{equation}
\boxed{K_{\lambda_n}=W_n0,\qquad
\lambda_n\nearrow\min\Lambda_\tau.}
\label{hmtfg:clasificacion:eq:10}
\end{equation}

\textbf{Démonstration.} Les lemmes de comparaison et d’existence placent chaque paramètre sélectionné au-dessous de \(a\), avec un itinéraire inférieur à \(\tau\). Le théorème de classification identifie cet itinéraire à \(W_n0\). La suite est croissante et majorée par \(a\) ; soit \(b\le a\) sa limite.

Fixons \(p\ge1\). Pour \(n\) suffisamment grand, \(2p<2^n\). Alors les signes de \(c_p(\lambda_n)\) et de \(c_{2p}(\lambda_n)\) coïncident avec \(\tau_p\) et \(\tau_{2p}=-\tau_p\), respectivement. La borne uniforme des dérivées secondes,
\[
\|(f_\lambda^p)''\|\le B_p,
\qquad B_p=16^p\frac{16^p-1}{15},
\]
et l’identité \((f_\lambda^p)'(0)=0\) impliquent
\[
|c_{2p}(\lambda_n)-c_p(\lambda_n)|
\le\tfrac12B_p|c_p(\lambda_n)|^2.
\]
Comme les deux termes du membre de gauche sont de signes opposés,
\begin{equation}
|c_p(\lambda_n)|>2/B_p.
\label{hmtfg:clasificacion:eq:11}
\end{equation}
Par continuité, \(c_p(b)\) conserve le signe \(\tau_p\) et satisfait \(|c_p(b)|\ge2/B_p>0\). Cela vaut pour chaque \(p\), de sorte que \(b\in\Lambda_\tau\). La minimalité de \(a\) donne \(b\ge a\), et donc \(b=a\). \(\square\)

La preuve complète conserve le sélecteur global et produit ses itinéraires et sa limite. L’information des huit premiers niveaux est compatible avec le résultat, mais la démonstration de \eqref{hmtfg:clasificacion:eq:10} n’extrapole pas ces huit niveaux : elle utilise la récurrence, des intervalles invariants et la séparation uniforme de chaque signe.


\subsection{Contrôles logiques et falsificateurs}
\label{hmtfg:clasificacion:7}

\par\noindent\textbf{1.}\enspace Le théorème de classification exige un maximum unimodal strict et un intervalle invariant. Si l’une des deux branches perd sa monotonie, les intervalles \eqref{hmtfg:clasificacion:eq:2} peuvent cesser d’être restrictifs et la classification doit être vérifiée à nouveau.
\par\noindent\textbf{2.}\enspace La condition \(K<\tau\) est essentielle. Au-delà du mot limite, il peut exister des centres de période \(2^n\) avec d’autres combinatoires ; le théorème ne les identifie pas à \(W_n0\).
\par\noindent\textbf{3.}\enspace L’analyticité en le paramètre assure l’isolement et la finitude des zéros sur le compact. Pour une famille seulement continue, l’existence d’une première racine strictement postérieure exige un argument supplémentaire.
\par\noindent\textbf{4.}\enspace Une limite de racines critiques peut, en général, perdre des signes en atteignant un retour de période inférieure. L’inégalité \eqref{hmtfg:clasificacion:eq:11} exclut précisément ce phénomène pour cette suite et chaque horizon fixé.
\par\noindent\textbf{5.}\enspace La convergence des premières racines vers le premier paramètre d’itinéraire \(\tau\) n’équivaut pas à l’unicité de tous les paramètres de cet itinéraire et ne détermine pas à elle seule un taux métrique. Le croisement transverse et l’identification de la branche apportent ces conclusions ultérieures.
